% ═══════════════════════════════════════════════════════════════
% LH-02a: Ohmsches Gesetz – LEHRERHINWEISE
% Physik Klasse 8 | Doppelstunde 02a (80 Min)
% ═══════════════════════════════════════════════════════════════

\documentclass[11pt, a4paper]{article}

\usepackage[utf8]{inputenc}
\usepackage[T1]{fontenc}
\usepackage[ngerman]{babel}

\usepackage{helvet}
\renewcommand{\familydefault}{\sfdefault}

\usepackage{microtype}
\usepackage[margin=2cm]{geometry}
\usepackage{enumitem}
\usepackage{tabularx}
\usepackage{booktabs}
\usepackage{array}
\usepackage{longtable}

\usepackage{amsmath}
\usepackage{siunitx}
\sisetup{locale = DE}

\usepackage{fancyhdr}
\usepackage{lastpage}

\usepackage{xcolor}
\usepackage{tcolorbox}
\tcbuselibrary{skins}

% Farben
\definecolor{physik}{HTML}{2c5aa0}
\definecolor{hellgrau}{HTML}{F5F5F5}
\definecolor{mittelgrau}{HTML}{888888}
\definecolor{success}{HTML}{2a7a4b}
\definecolor{warning}{HTML}{b35c00}
\definecolor{error}{HTML}{c62828}

\colorlet{fachfarbe}{physik}

% Kopfzeile
\pagestyle{fancy}
\fancyhf{}
\renewcommand{\headrulewidth}{0.4pt}
\lhead{\small Physik Kl. 8 | Das Ohmsche Gesetz | \textbf{Lehrerhinweise}}
\rhead{\small 80 Min | Doppelstunde 02a}
\cfoot{\small Seite \thepage{} von \pageref{LastPage}}

% Boxen
\newtcolorbox{infobox}[1][]{
    colback=hellgrau,
    colframe=hellgrau,
    boxrule=0pt,
    arc=0pt,
    left=8pt, right=8pt, top=6pt, bottom=6pt,
    borderline west={3pt}{0pt}{fachfarbe},
    #1
}

\newtcolorbox{warnbox}[1][]{
    colback=white,
    colframe=warning,
    boxrule=1pt,
    arc=0pt,
    left=8pt, right=8pt, top=6pt, bottom=6pt,
    #1
}

\newtcolorbox{tipbox}[1][]{
    colback=white,
    colframe=success,
    boxrule=1pt,
    arc=0pt,
    left=8pt, right=8pt, top=6pt, bottom=6pt,
    #1
}

\setlength{\parindent}{0pt}
\setlength{\parskip}{6pt}

\begin{document}

% ═══════════════════════════════════════════════════════════════
% HEADER
% ═══════════════════════════════════════════════════════════════

\begin{center}
{\Large\bfseries\textcolor{fachfarbe}{Lehrerhinweise: Das Ohmsche Gesetz}}\\[0.3em]
{\small Physik Klasse 8 | Doppelstunde 02a | 80 Minuten}
\end{center}
\vspace{0.3em}
\hrule
\vspace{1em}

% ═══════════════════════════════════════════════════════════════
% ÜBERBLICK
% ═══════════════════════════════════════════════════════════════

\section*{Überblick}

\begin{tabularx}{\textwidth}{lX}
\textbf{Thema:} & Das Ohmsche Gesetz, I(U)-Kennlinien \\
\textbf{Vorwissen:} & Spannung, Stromstärke, Reihen-/Parallelschaltung \\
\textbf{Lernziele:} & Proportionalität I $\sim$ U, Formel R = U/I, Kennlinien zeichnen \\
\textbf{Differenzierung:} & $\star$ 26 BE | $\star\star$ 39 BE | $\star\star\star$ 45 BE \\
\end{tabularx}

% ═══════════════════════════════════════════════════════════════
% MATERIAL
% ═══════════════════════════════════════════════════════════════

\section*{Benötigtes Material}

\begin{itemize}[leftmargin=1.5em, itemsep=2pt]
    \item \textbf{Digital:} Lernpfad LP-02a (auf Schüler-Tablets/PCs)
    \item \textbf{Print:} Arbeitsblatt AB-02a (1× pro SuS)
    \item \textbf{Stifte:} Bleistift, Lineal, 2 verschiedenfarbige Stifte
    \item \textbf{Optional:} Beamer für gemeinsame Besprechung
\end{itemize}

\begin{tipbox}
\textbf{Tipp:} Simulationen vorab testen! URLs:\\
\texttt{SIM-02a-kennlinie.html} und \texttt{SIM-02a-gluehlampe.html}
\end{tipbox}

% ═══════════════════════════════════════════════════════════════
% STUNDENVERLAUF
% ═══════════════════════════════════════════════════════════════

\section*{Stundenverlauf (80 Min)}

\renewcommand{\arraystretch}{1.3}
\begin{longtable}{|c|p{3cm}|p{7cm}|p{3cm}|}
\hline
\textbf{Zeit} & \textbf{Phase} & \textbf{Inhalt} & \textbf{Material} \\
\hline
\endfirsthead
\hline
\textbf{Zeit} & \textbf{Phase} & \textbf{Inhalt} & \textbf{Material} \\
\hline
\endhead
0--5 & Einstieg & Taschenlampen-Frage (LP Slide 1). Warum leuchtet eine Lampe mit neuen Batterien heller? & LP Slide 1 \\
\hline
5--12 & Erarbeitung I & Wassermodell einführen (LP Slides 2--3). Analogie: Druck $\rightarrow$ U, Durchfluss $\rightarrow$ I, Engstelle $\rightarrow$ R & LP + AB Kasten 1 \\
\hline
12--25 & Experiment & Simulation SIM-02a-kennlinie: 6 Messwerte aufnehmen (LP Slide 4) & LP + AB Aufg. 1 \\
\hline
25--35 & Auswertung & Kennlinie zeichnen (LP Slide 5). Ausgleichsgerade mit Lineal! & AB Aufg. 2 \\
\hline
35--40 & \textit{Pause} & 5 Min Pause & --- \\
\hline
40--50 & Erarbeitung II & Ohmsches Gesetz formulieren (LP Slides 6--7). Alle 3 Formeln! & LP + AB Kasten 2 \\
\hline
50--60 & Übung & Berechnungen: Aufg. 3a--d gemeinsam, dann selbstständig (LP Slides 8--9) & AB Aufg. 3 \\
\hline
60--70 & Erweiterung & Glühlampe: SIM-02a-gluehlampe, 2. Kennlinie zeichnen (LP Slides 10--12) & LP + AB Aufg. 4 \\
\hline
70--78 & Sicherung & Erklärung Atomrümpfe, Lückentext ausfüllen (LP Slide 13) & AB Aufg. 4c \\
\hline
78--80 & Abschluss & Zusammenfassung (LP Slide 14), AB Kasten 3 ergänzen & AB Kasten 3 \\
\hline
\end{longtable}

% ═══════════════════════════════════════════════════════════════
% DIFFERENZIERUNG
% ═══════════════════════════════════════════════════════════════

\section*{Differenzierung}

\begin{tabularx}{\textwidth}{|c|X|c|}
\hline
\textbf{Niveau} & \textbf{Aufgaben} & \textbf{BE} \\
\hline
$\star$ Basis & Aufg. 1, 2, 3a--d + Kästen 1--3 & 26 BE \\
\hline
$\star\star$ Standard & Zusätzlich: Aufg. 3e--g, Aufg. 4 komplett & 39 BE \\
\hline
$\star\star\star$ Erweitert & Zusätzlich: Aufg. 5 (Transfer) & 45 BE \\
\hline
\end{tabularx}

\begin{infobox}
\textbf{Hinweis:} Bei gemischten Kursen können $\star$-SuS nach Aufg. 2 die Glühlampen-Simulation überspringen und direkt zu Kasten 3 gehen.
\end{infobox}

\begin{tipbox}
\textbf{Formelumstellung – Differenzierung:}
\begin{itemize}[leftmargin=1.5em, itemsep=2pt]
    \item $\star$ \textbf{Basis:} Das Formeldreieck wird \textit{nahegelegt}. SuS nutzen es als Merkhilfe.
    \item $\star\star$ \textbf{Standard:} Das Formeldreieck ist \textit{optional}. SuS dürfen wählen.
    \item $\star\star\star$ \textbf{Erweitert:} Formale Umstellung wird \textit{verlangt}. Rechenweg zeigen!
\end{itemize}
\textit{Alle Niveaus müssen die Formeln anwenden können – nur die Methode unterscheidet sich.}
\end{tipbox}

% ═══════════════════════════════════════════════════════════════
% HÄUFIGE FEHLER
% ═══════════════════════════════════════════════════════════════

\section*{Häufige Fehler und Reaktionen}

\begin{warnbox}
\textbf{Fehler 1: Formel vertauscht}\\
SuS rechnen $R = I/U$ statt $R = U/I$\\
\textit{Reaktion:} Formeldreieck zeigen. "`Was suchst du? Das deckst du ab!"'
\end{warnbox}

\begin{warnbox}
\textbf{Fehler 2: Einheiten vergessen}\\
SuS schreiben nur Zahlen, keine Einheiten.\\
\textit{Reaktion:} "`Die Einheit gehört zum Ergebnis wie der Nachname zum Vornamen."'
\end{warnbox}

\begin{warnbox}
\textbf{Fehler 3: Kennlinie nicht durch Ursprung}\\
SuS zeichnen Gerade, die nicht bei (0,0) beginnt.\\
\textit{Reaktion:} "`Was passiert, wenn U = 0? Dann ist I = ..."'
\end{warnbox}

\begin{warnbox}
\textbf{Fehler 4: Glühlampen-Kennlinie als Gerade}\\
SuS verbinden Punkte geradlinig statt als Kurve.\\
\textit{Reaktion:} "`Schau dir die Punkte genau an. Liegen sie auf einer Geraden?"'
\end{warnbox}

% ═══════════════════════════════════════════════════════════════
% LÖSUNGSÜBERSICHT
% ═══════════════════════════════════════════════════════════════

\section*{Lösungsübersicht}

\subsection*{Kasten 1 (Wassermodell)}
\begin{itemize}[leftmargin=1.5em]
    \item Wasserdruck $\rightarrow$ \textbf{Spannung U}
    \item Wassermenge/s $\rightarrow$ \textbf{Stromstärke I}
    \item Engstelle $\rightarrow$ \textbf{Widerstand R}
\end{itemize}

\subsection*{Aufgabe 1 (Wertetabelle für R = 10 $\Omega$)}
\begin{center}
\begin{tabular}{|c|c|c|c|c|c|c|}
\hline
U (V) & 2 & 4 & 6 & 8 & 10 & 12 \\
\hline
I (A) & 0,2 & 0,4 & 0,6 & 0,8 & 1,0 & 1,2 \\
\hline
R ($\Omega$) & 10 & 10 & 10 & 10 & 10 & 10 \\
\hline
\end{tabular}
\end{center}

\subsection*{Aufgabe 3}
\begin{itemize}[leftmargin=1.5em]
    \item a) $R = U/I$, $U = R \cdot I$, $I = U/R$
    \item b) $R = 12V / 4A = 3\,\Omega$
    \item c) $U = 6\Omega \cdot 2A = 12\,$V
    \item d) $I = 24V / 8\Omega = 3\,$A
    \item e) $U = 50\Omega \cdot 0{,}4A = 20\,$V
    \item f) $R = 10V / 0{,}5A = 20\,\Omega$
    \item g) $I = 20V / 100\Omega = 0{,}2A = 200\,$mA
\end{itemize}

\subsection*{Aufgabe 4 (Glühlampe)}
\begin{center}
\begin{tabular}{|c|c|c|c|c|c|c|}
\hline
U (V) & 2 & 4 & 6 & 8 & 10 & 12 \\
\hline
I (A) & 0,35 & 0,55 & 0,68 & 0,78 & 0,85 & 0,90 \\
\hline
\end{tabular}
\end{center}

\textbf{Lückentext:} Strom, Atomrümpfen, schwingen, heiß, größer, weniger

\subsection*{Aufgabe 5}
\begin{itemize}[leftmargin=1.5em]
    \item a) Föhn: $R = 230V / 10A = 23\,\Omega$
    \item b) Lisa hat U und I vertauscht oder sich verrechnet
    \item c) Kabel hat Eigenwiderstand, bei hohem Strom entsteht Wärme
\end{itemize}

% ═══════════════════════════════════════════════════════════════
% NACHBEREITUNG
% ═══════════════════════════════════════════════════════════════

\section*{Nachbereitung / Ausblick}

\begin{itemize}[leftmargin=1.5em]
    \item \textbf{Nächste Stunde:} Widerstandsmessung, spezifischer Widerstand
    \item \textbf{Hausaufgabe:} AB Aufg. 5 (für die, die nicht fertig wurden)
    \item \textbf{Vertiefung:} PhET-Simulation "`Ohm's Law"' für Zusatzübung
\end{itemize}

\vfill
\hrule
\vspace{0.3em}
{\small\textcolor{mittelgrau}{LH-02a | Stand: \today}}

\end{document}
